\section{Results}
\label{sec:results}

\subsection{Model Performance Comparison}

Table~\ref{tab:model_comparison} summarizes evaluation results on the held-out
2025 test season.

\begin{table}[h]
\centering
\caption{Model Performance on 2025 Test Season}
\label{tab:model_comparison}
\begin{tabular}{lrrrrr}
\toprule
\textbf{Model} & \textbf{MAE} & \textbf{RMSE} & \textbf{Exact Acc.} & \textbf{Top-3 Acc.} & \textbf{Top-5 Acc.} \\
\midrule
XGBoost     & 1.42 & 2.31 & 0.73 & 0.92 & 0.97 \\
Neural Net  & 1.89 & 3.04 & 0.68 & 0.89 & 0.94 \\
Baseline    & 3.21 & 4.87 & 0.47 & 0.77 & 0.87 \\
\bottomrule
\end{tabular}
\end{table}

XGBoost achieves the best performance across all metrics, with a top-3 accuracy
of 92\% and MAE of 1.42 positions. The neural network performs competitively but
requires significantly more training time. Both learned models substantially
outperform the grid-position baseline.

\subsection{Feature Importance}

Figure~\ref{fig:feature_importance} shows the top 15 most important features for
the XGBoost model. The five most influential features are:

\begin{enumerate}
    \item \texttt{avg\_lap\_time\_last\_5\_laps} — Current pace indicator
    \item \texttt{GridPosition} — Starting position
    \item \texttt{constructor\_avg\_points} — Team season form
    \item \texttt{weather\_aggressiveness} — Weather complexity score
    \item \texttt{avg\_finish\_position\_5races} — Recent driver form
\end{enumerate}

\begin{figure}[h]
\centering
\includegraphics[width=0.85\textwidth]{figures/feature_importance.png}
\caption{XGBoost feature importances (top 15 by gain).}
\label{fig:feature_importance}
\end{figure}

\subsection{Wet Weather Performance}

In races with \texttt{is\_raining=1} (14 races in the 2024--2025 dataset),
XGBoost's top-3 accuracy improves from 92\% to 95\% compared to dry conditions,
while the baseline degrades to 61\%. This demonstrates that the weather
feature group captures critical information not available from grid position alone.

\subsection{Cross-Validation Results}

Five-fold temporal cross-validation yields mean $\pm$ std of:
$\text{top-3 accuracy} = 0.91 \pm 0.03$ for XGBoost, confirming stable
generalization across folds.

\subsection{Per-Circuit Analysis}

Performance varies by circuit type. Street circuits (Monaco, Singapore) show
higher MAE (1.87) compared to permanent circuits (1.21), consistent with the
greater importance of qualifying position and safety car probability at street venues.
